% GENERATED FILE. Edit canonical lessons, then run npm run generate:books.
% canonical-source-hash: fnv1a64:f3f6621086a04eb1
% canonical-lessons: GE-C480-uebelkeit, GE-C480-uebergeben, GE-C480-durchfall, GE-C480-schwindlig, GE-C480-kreislauf

\chapter{Übel und schwindlig}
\label{ch:ge-ubel-und-schwindlig}
\hlchaptermodality{480}

\begin{chapteropening}
\textbf{By the end of this chapter:} \emph{I can describe feeling sick and dizzy.}

Complete the last lesson of chapter 480: I can describe feeling sick and dizzy.
\end{chapteropening}

\section[die Übelkeit]{die Übelkeit — nausea}

\label{lesson:GE-C480-uebelkeit}



\begin{quote}
Before the new one: say the German for sweat, then the German for goose bumps.
\end{quote}

\subsection*{You'll want to know: die Übelkeit}
\textbf{die Übelkeit} — \textquotedblleft{}nausea\textquotedblright{}.

\begin{quote}
Auf dem Schiff hatte ich starke Übelkeit.
\end{quote}

\begin{quote}
\emph{On the ship I felt very sick.}
\end{quote}

No plural. \emph{Mir ist übel} — I feel sick.

\begin{cousinweb}
\emph{Übelkeit} is built on \emph{übel} \textquotedblleft{}bad, sick\textquotedblright{} with \emph{-keit}; \emph{übel} is inherited from Old High German \emph{ubil}, from Proto-Germanic \emph{ubilaz} (EtymDB), and English \emph{evil} is a cousin.
\end{cousinweb}

\begin{grammarlens}[title={What you've built}]
One more word for telling, explaining and describing at B1.
\end{grammarlens}

\subsection*{Your turn}
\begin{itemize}
\raggedright
  \item \emph{Say it:} \emph{die Übelkeit}
  \item \emph{Say it:} \emph{die Übelkeit}, once more
  \item \emph{Say it:} a sentence of your own with \emph{die Übelkeit}
  \item \emph{Recall:} say the German for sweat, then the German for goose bumps, then say \emph{die Übelkeit} again
\end{itemize}

\begin{tcolorbox}[breakable,colback=teal!4,colframe=teal!35!black,title={Before you move on}]
What does die Übelkeit mean? (\textquotedblleft{}Nausea\textquotedblright{}.) Say it once more.
\end{tcolorbox}
\section[übergeben]{übergeben — to be sick, to vomit (sich übergeben); to hand over}

\label{lesson:GE-C480-uebergeben}



\begin{quote}
Before the new one: say the German for goose bumps, then the German for nausea.
\end{quote}

\subsection*{You'll want to know: übergeben}
\textbf{übergeben} — \textquotedblleft{}to be sick, to vomit (sich übergeben); to hand over\textquotedblright{}.

\begin{quote}
Im Bus hat sich ein Kind übergeben.
\end{quote}

\begin{quote}
\emph{A child was sick on the bus.}
\end{quote}

Inseparable. Reflexive for to vomit: \emph{Er hat sich übergeben.} Without \emph{sich} it means to hand over: \emph{den Schlüssel übergeben}.

\begin{cousinweb}
\emph{übergeben} is built on \emph{über} \textquotedblleft{}over\textquotedblright{} and \emph{geben} \textquotedblleft{}to give\textquotedblright{} (EtymDB), inherited from Old High German \emph{geban}, from Proto-Germanic \emph{gebaną}.
\end{cousinweb}

\begin{grammarlens}[title={What you've built}]
One more word for telling, explaining and describing at B1.
\end{grammarlens}

\subsection*{Your turn}
\begin{itemize}
\raggedright
  \item \emph{Say it:} \emph{übergeben}
  \item \emph{Say it:} \emph{übergeben}, once more
  \item \emph{Say it:} a sentence of your own with \emph{übergeben}
  \item \emph{Recall:} say the German for goose bumps, then the German for nausea, then say \emph{übergeben} again
\end{itemize}

\begin{tcolorbox}[breakable,colback=teal!4,colframe=teal!35!black,title={Before you move on}]
What does übergeben mean? (\textquotedblleft{}To be sick, to vomit (sich übergeben); to hand over\textquotedblright{}.) Say it once more.
\end{tcolorbox}
\section[der Durchfall]{der Durchfall — diarrhoea}

\label{lesson:GE-C480-durchfall}



\begin{quote}
Before the new one: say the German for nausea, then the German for to be sick.
\end{quote}

\subsection*{You'll want to know: der Durchfall}
\textbf{der Durchfall} — \textquotedblleft{}diarrhoea\textquotedblright{}.

\begin{quote}
Nach dem Essen im Urlaub hatte er Durchfall.
\end{quote}

\begin{quote}
\emph{After the food on holiday he had diarrhoea.}
\end{quote}

No article in \emph{Durchfall haben}. No plural in everyday use.

\begin{cousinweb}
\emph{Durchfall} is built on \emph{durch} \textquotedblleft{}through\textquotedblright{} and \emph{Fall} \textquotedblleft{}fall\textquotedblright{}, which is inherited from Middle High German \emph{val} and Old High German \emph{val} (EtymDB).
\end{cousinweb}

\begin{grammarlens}[title={What you've built}]
One more word for telling, explaining and describing at B1.
\end{grammarlens}

\subsection*{Your turn}
\begin{itemize}
\raggedright
  \item \emph{Say it:} \emph{der Durchfall}
  \item \emph{Say it:} \emph{der Durchfall}, once more
  \item \emph{Say it:} a sentence of your own with \emph{der Durchfall}
  \item \emph{Recall:} say the German for nausea, then the German for to be sick, then say \emph{der Durchfall} again
\end{itemize}

\begin{tcolorbox}[breakable,colback=teal!4,colframe=teal!35!black,title={Before you move on}]
What does der Durchfall mean? (\textquotedblleft{}Diarrhoea\textquotedblright{}.) Say it once more.
\end{tcolorbox}
\section[schwindlig]{schwindlig — dizzy}

\label{lesson:GE-C480-schwindlig}



\begin{quote}
Before the new one: say the German for to be sick, then the German for diarrhoea.
\end{quote}

\subsection*{You'll want to know: schwindlig}
\textbf{schwindlig} — \textquotedblleft{}dizzy\textquotedblright{}.

\begin{quote}
Wenn ich zu schnell aufstehe, wird mir schwindlig.
\end{quote}

\begin{quote}
\emph{When I stand up too fast, I get dizzy.}
\end{quote}

With a dative: \emph{Mir ist schwindlig.} Also spelt \emph{schwindelig}.

\begin{cousinweb}
\emph{schwindlig} is built on \emph{Schwindel} \textquotedblleft{}giddiness\textquotedblright{} with \emph{-ig}; \emph{Schwindel} comes from \emph{schwindeln}, inherited from Middle High German \emph{swindeln} and Old High German \emph{swintilōn} (EtymDB).
\end{cousinweb}

\begin{grammarlens}[title={What you've built}]
One more word for telling, explaining and describing at B1.
\end{grammarlens}

\subsection*{Your turn}
\begin{itemize}
\raggedright
  \item \emph{Say it:} \emph{schwindlig}
  \item \emph{Say it:} \emph{schwindlig}, once more
  \item \emph{Say it:} a sentence of your own with \emph{schwindlig}
  \item \emph{Recall:} say the German for to be sick, then the German for diarrhoea, then say \emph{schwindlig} again
\end{itemize}

\begin{tcolorbox}[breakable,colback=teal!4,colframe=teal!35!black,title={Before you move on}]
What does schwindlig mean? (\textquotedblleft{}Dizzy\textquotedblright{}.) Say it once more.
\end{tcolorbox}
\section[der Kreislauf]{der Kreislauf — circulation; cycle}

\label{lesson:GE-C480-kreislauf}



\begin{quote}
Before the new one: say the German for diarrhoea, then the German for dizzy.
\end{quote}

\subsection*{You'll want to know: der Kreislauf}
\textbf{der Kreislauf} — \textquotedblleft{}circulation; cycle\textquotedblright{}.

\begin{quote}
Bei der Hitze hatte meine Oma Probleme mit dem Kreislauf.
\end{quote}

\begin{quote}
\emph{In the heat my grandma had problems with her circulation.}
\end{quote}

Plural: \emph{die Kreisläufe}. Germans often blame \emph{der Kreislauf} when they feel faint.

\begin{cousinweb}
\emph{Kreislauf} is built on \emph{Kreis} \textquotedblleft{}circle\textquotedblright{}, inherited from Old High German \emph{kreiz}, and \emph{Lauf} \textquotedblleft{}run\textquotedblright{}, from \emph{laufen}, from Proto-Germanic \emph{hlaupaną} (EtymDB).
\end{cousinweb}

\begin{grammarlens}[title={What you've built}]
One more word for telling, explaining and describing at B1.
\end{grammarlens}

\subsection*{Your turn}
\begin{itemize}
\raggedright
  \item \emph{Say it:} \emph{der Kreislauf}
  \item \emph{Say it:} \emph{der Kreislauf}, once more
  \item \emph{Say it:} a sentence of your own with \emph{der Kreislauf}
  \item \emph{Recall:} say the German for diarrhoea, then the German for dizzy, then say \emph{der Kreislauf} again
\end{itemize}

\begin{tcolorbox}[breakable,colback=teal!4,colframe=teal!35!black,title={Before you move on}]
What does der Kreislauf mean? (\textquotedblleft{}Circulation; cycle\textquotedblright{}.) Say it once more.
\end{tcolorbox}
